\documentclass[12pt]{article}
\usepackage[utf8]{inputenc}
\usepackage[margin=1in]{geometry}
\usepackage{amsmath, amssymb, amsthm}
\usepackage{xcolor}
\usepackage[most]{tcolorbox}
\usepackage{enumitem}
\usepackage{fancyhdr}
\usepackage{titlesec}

% Define colored boxes
\newtcolorbox{problemconfigbox}{
    colback=blue!5!white,
    colframe=blue!75!black,
    title=\textbf{1. PROBLEM CONFIGURATION},
    fonttitle=\bfseries,
    breakable
}

\newtcolorbox{strategicbox}{
    colback=pink!5!white,
    colframe=pink!75!black,
    title=\textbf{2. STRATEGIC FRAMEWORK APPLICATION},
    fonttitle=\bfseries,
    breakable
}

\newtcolorbox{tacticalbox}{
    colback=green!5!white,
    colframe=green!75!black,
    title=\textbf{3. TACTICAL METHODOLOGY SELECTION},
    fonttitle=\bfseries,
    breakable
}

\newtcolorbox{conversionbox}{
    colback=yellow!5!white,
    colframe=yellow!75!black,
    title=\textbf{4. CONVERSION TECHNIQUES APPLICATION},
    fonttitle=\bfseries,
    breakable
}

\newtcolorbox{verificationbox}{
    colback=purple!5!white,
    colframe=purple!75!black,
    title=\textbf{5. VERIFICATION STRATEGY},
    fonttitle=\bfseries,
    breakable
}

\newtcolorbox{roadmapbox}{
    colback=orange!5!white,
    colframe=orange!75!black,
    title=\textbf{6. SOLUTION ROADMAP},
    fonttitle=\bfseries,
    breakable
}

\newtcolorbox{competitionbox}{
    colback=cyan!5!white,
    colframe=cyan!75!black,
    title=\textbf{7. COMPETITION INTEGRATION},
    fonttitle=\bfseries,
    breakable
}

\newtcolorbox{highlightbox}[1]{
    colback=red!5!white,
    colframe=red!75!black,
    title=#1,
    fonttitle=\bfseries
}

\title{\textbf{Strategic Analysis: 2018 AMC 12A Problem 14\\Logarithmic Equation}}
\author{Generated using AMC12/COMC Problem Solving Framework}
\date{\today}

\begin{document}

\maketitle

\section*{Problem Statement}

The solution to the equation $\log_{3x} 4 = \log_{2x} 8$, where $x$ is a positive real number other than $\frac{1}{3}$ or $\frac{1}{2}$, can be written as $\frac{p}{q}$ where $p$ and $q$ are relatively prime positive integers. What is $p + q$?

\textbf{Answer Choices:}
\begin{itemize}
    \item[(A)] $5$
    \item[(B)] $13$
    \item[(C)] $17$
    \item[(D)] $31$
    \item[(E)] $35$
\end{itemize}

\begin{problemconfigbox}
\subsection*{A. Problem Classification}
\begin{itemize}[leftmargin=*]
    \item \textbf{Problem Type}: To Find (solve for $x$, then compute $p+q$)
    \item \textbf{Difficulty Band}: Mid-band (Problem 14 out of 25)
    \item \textbf{Competition Context}: AMC 12A 2018
    \item \textbf{Mathematical Domain}: Algebra - Logarithmic Equations with emphasis on logarithm properties and change of base formula
\end{itemize}

\subsection*{B. Problem Structure Identification}
\begin{itemize}[leftmargin=*]
    \item \textbf{Given Information}: 
    \begin{itemize}
        \item Equation: $\log_{3x} 4 = \log_{2x} 8$
        \item Domain restriction: $x > 0$, $x \neq \frac{1}{3}$, $x \neq \frac{1}{2}$
        \item Solution is rational: $x = \frac{p}{q}$ where $\gcd(p,q) = 1$
    \end{itemize}
    \item \textbf{Constraints}: 
    \begin{itemize}
        \item Bases must be positive and $\neq 1$: $3x > 0$, $3x \neq 1$, $2x > 0$, $2x \neq 1$
        \item Solution must be in lowest terms
    \end{itemize}
    \item \textbf{Objective}: Find $p + q$ where $x = \frac{p}{q}$ in lowest terms
    \item \textbf{Hidden Assumptions}: 
    \begin{itemize}
        \item The equation has a unique solution (implied by ``the solution'')
        \item We need to simplify $x$ to lowest terms before computing $p+q$
        \item The given restrictions ensure bases are valid
    \end{itemize}
    \item \textbf{Answer Choice Leverage}: 
    \begin{itemize}
        \item All answers are relatively small positive integers
        \item (D) $31 = 4 + 27$ suggests $x = \frac{4}{27}$
        \item (E) $35 = 8 + 27$ or $5 + 30$ are alternatives
        \item Can work backwards to verify
    \end{itemize}
\end{itemize}

\subsection*{C. Strategic Orientation}
\begin{itemize}[leftmargin=*]
    \item \textbf{Hypothesis}: Use change of base formula to convert to common base
    \item \textbf{Conclusion}: Solve resulting equation for $x$, express as fraction, find sum
    \item \textbf{Notation}: 
    \begin{itemize}
        \item Let $a = 3x$, $b = 2x$ (bases)
        \item Note $4 = 2^2$, $8 = 2^3$ (powers of 2)
    \end{itemize}
    \item \textbf{Intuitive Approach}: 
    \begin{itemize}
        \item Recognize $4$ and $8$ are powers of $2$
        \item Apply change of base to common base
        \item Use logarithm identities to simplify
    \end{itemize}
    \item \textbf{Cross-Domain Potential}: 
    \begin{itemize}
        \item Pure algebraic manipulation
        \item Exponential form approach
        \item Change of base to natural log or $\log_2$
    \end{itemize}
\end{itemize}
\end{problemconfigbox}

\begin{strategicbox}
\subsection*{A. Psychological Strategies}
\begin{itemize}[leftmargin=*]
    \item \textbf{Mental Toughness}: Don't be intimidated by unusual bases ($3x$ and $2x$). These are just expressions, not standard bases.
    \item \textbf{Creativity Cultivation}: Notice that $4 = 2^2$ and $8 = 2^3$ - this suggests rewriting in terms of powers of 2
    \item \textbf{Iterative Refinement}: If change of base seems messy, try exponential form approach
\end{itemize}

\subsection*{B. Investigation Strategies}
\begin{itemize}[leftmargin=*]
    \item \textbf{Startup Orientation}: Identify logarithm equation with variable bases - key is change of base formula
    \item \textbf{Get Your Hands Dirty}: Write $\log_{3x} 4 = \frac{\log 4}{\log(3x)}$ immediately
    \item \textbf{Penultimate Step}: Final step before answer is simplifying $\frac{4}{27}$ and computing $4+27$
    \item \textbf{Wishful Thinking}: ``I wish the bases were the same'' - change of base makes them comparable!
    \item \textbf{Make It Easier}: Rewrite $4 = 2^2$ and $8 = 2^3$ to use power rule
    \item \textbf{Conjecture via Small Cases}: Try $x=1$: $\log_3 4 \approx 1.26$, $\log_2 8 = 3$, not equal
    \item \textbf{Visualization}: Not particularly helpful for algebraic equation
\end{itemize}

\subsection*{C. Late-Band Strategic Playbook}
Not applicable - this is Problem 14 (mid-band), not 17-25.

\subsection*{D. Argument Construction Strategy}
\begin{itemize}[leftmargin=*]
    \item \textbf{Direct Proof}: Apply change of base, simplify algebraically to $x = \frac{4}{27}$
    \item \textbf{Proof by Contradiction}: Not needed
    \item \textbf{Mathematical Induction}: Not applicable
    \item \textbf{Stylistic Conventions}: Use change of base consistently throughout
\end{itemize}

\subsection*{E. Cross-Domain Translation}
\begin{itemize}[leftmargin=*]
    \item \textbf{Draw a Picture}: Not helpful for this problem
    \item \textbf{Recast the Problem}: ``Find $x$ such that $(3x)^a = 4$ and $(2x)^a = 8$ for some $a$''
    \item \textbf{Change Your Point of View}: 
    \begin{itemize}
        \item View as exponential equation via $y = \log_{3x} 4 = \log_{2x} 8$
        \item View as relationship between two logarithms
        \item View using change of base to common base
    \end{itemize}
\end{itemize}
\end{strategicbox}

\begin{tacticalbox}
\subsection*{A. Core Mathematical Tactics}
\begin{itemize}[leftmargin=*]
    \item \textbf{Symmetry}: Not directly applicable
    \item \textbf{The Extreme Principle}: Not applicable
    \item \textbf{The Pigeonhole Principle}: Not applicable
    \item \textbf{Invariants}: The equality relationship is preserved through change of base
\end{itemize}

\subsection*{B. Crossover Tactics}
\begin{itemize}[leftmargin=*]
    \item \textbf{Graph Theory}: Not applicable
    \item \textbf{Complex Numbers}: Not applicable
    \item \textbf{Generating Functions}: Not applicable
\end{itemize}
\end{tacticalbox}

\begin{conversionbox}
\subsection*{A. Recognition Index}

\textbf{High-Probability Techniques Applicable:}
\begin{enumerate}
    \item \textbf{Logarithm Identities} (100\% match): Core technique for this problem
    \item \textbf{Change of Base Formula} (100\% match): Essential tool
    \item \textbf{Power Rule for Logarithms} (95\% match): $\log_b(a^n) = n\log_b a$
    \item \textbf{Algebraic Manipulation} (90\% match): Simplifying resulting equation
    \item \textbf{Exponential Form Conversion} (85\% match): Alternative approach
\end{enumerate}

\subsection*{B. Technique Selection}

\textbf{Primary Approach: Change of Base Formula}
\begin{itemize}
    \item Success Rate: 98\%
    \item Time Cost: 2-3 minutes
    \item Cognitive Load: Medium (requires careful algebra)
    \item Justification: Most direct and systematic approach
\end{itemize}

\textbf{Secondary Approach: Power Manipulation with Equal Bases}
\begin{itemize}
    \item Success Rate: 95\%
    \item Time Cost: 2-3 minutes
    \item Cognitive Load: Medium (elegant insight required)
    \item Justification: Uses $\log_{b^n}(a^n) = \log_b a$ identity
\end{itemize}

\textbf{Tertiary Approach: Exponential Form}
\begin{itemize}
    \item Success Rate: 90\%
    \item Time Cost: 3-4 minutes
    \item Cognitive Load: Medium-High (requires careful manipulation)
    \item Justification: Alternative perspective, good for verification
\end{itemize}

\subsection*{C. Integration Strategy}

\textbf{Technique Chain (Primary Method):}
\begin{enumerate}
    \item Apply change of base: $\frac{\log_2 4}{\log_2(3x)} = \frac{\log_2 8}{\log_2(2x)}$
    \item Simplify using $\log_2 4 = 2$ and $\log_2 8 = 3$
    \item Cross-multiply: $3\log_2(3x) = 2\log_2(2x)$
    \item Apply power rule: $\log_2[(3x)^3] = \log_2[(2x)^2]$
    \item Set arguments equal: $(3x)^3 = (2x)^2$
    \item Expand: $27x^3 = 4x^2$
    \item Solve: $27x = 4$, so $x = \frac{4}{27}$
    \item Answer: $p + q = 4 + 27 = 31$
\end{enumerate}

\textbf{Technique Chain (Power Manipulation Method):}
\begin{enumerate}
    \item Rewrite $4 = 2^2$ and $8 = 2^3$: $\log_{3x} 2^2 = \log_{2x} 2^3$
    \item Apply power rule: $2\log_{3x} 2 = 3\log_{2x} 2$
    \item Recognize this relates the bases through powers
    \item Cube first equation's base, square second: $\log_{(3x)^3} 2^6 = \log_{(2x)^2} 2^6$
    \item Therefore $(3x)^3 = (2x)^2$
    \item Solve as before: $x = \frac{4}{27}$
\end{enumerate}

\subsection*{D. Conflict Resolution}

\textbf{Potential Conflicts:}
\begin{itemize}
    \item \textbf{Conflict}: Which base to use for change of base?
    \item \textbf{Resolution}: Base 2 is most natural since $4=2^2$ and $8=2^3$
    \item \textbf{Conflict}: Should I expand $(3x)^3$ and $(2x)^2$ immediately?
    \item \textbf{Resolution}: Yes, this gives $27x^3 = 4x^2$, easy to solve
\end{itemize}
\end{conversionbox}

\begin{verificationbox}
\subsection*{A. Level Selection with Decision Tree}

\textbf{Confidence Level: High (mid-band problem with standard technique)}

\textbf{Verification Level: Standard (Level 2)}
\begin{itemize}
    \item Check algebraic manipulations
    \item Verify $x = \frac{4}{27}$ satisfies original equation
    \item Confirm $\gcd(4, 27) = 1$
    \item Verify domain restrictions satisfied
\end{itemize}

\subsection*{B. Specialized Verification Methods}

\textbf{Direct Substitution:}
\begin{align*}
\text{LHS: } \log_{3 \cdot \frac{4}{27}} 4 &= \log_{\frac{4}{9}} 4 = \frac{\log_2 4}{\log_2 \frac{4}{9}} = \frac{2}{\log_2 4 - \log_2 9} \\
&= \frac{2}{2 - \log_2 9}\\
\text{RHS: } \log_{2 \cdot \frac{4}{27}} 8 &= \log_{\frac{8}{27}} 8 = \frac{\log_2 8}{\log_2 \frac{8}{27}} = \frac{3}{\log_2 8 - \log_2 27}\\
&= \frac{3}{3 - \log_2 27}
\end{align*}
We need to verify these are equal. Let $a = \log_2 3$. Then:
\begin{align*}
\text{LHS} &= \frac{2}{2 - 2a}\\
\text{RHS} &= \frac{3}{3 - 3a}
\end{align*}
These are equal: $\frac{2}{2-2a} = \frac{1}{1-a}$ and $\frac{3}{3-3a} = \frac{1}{1-a}$ $\checkmark$

\textbf{Domain Check:}
\begin{itemize}
    \item $x = \frac{4}{27} > 0$ $\checkmark$
    \item $3x = \frac{4}{9} \neq 1$ $\checkmark$
    \item $2x = \frac{8}{27} \neq 1$ $\checkmark$
    \item $x = \frac{4}{27} \neq \frac{1}{3}, \frac{1}{2}$ $\checkmark$
\end{itemize}

\textbf{Answer Form Check:}
\begin{itemize}
    \item $\gcd(4, 27) = 1$ $\checkmark$ (4 = $2^2$, 27 = $3^3$, no common factors)
    \item $p + q = 4 + 27 = 31$ is among choices (D) $\checkmark$
\end{itemize}

\subsection*{C. Confidence Calibration}

\textbf{Pre-Solution Confidence:} 75\% (standard logarithm problem)

\textbf{Post-Solution Confidence:} 98\% (algebra clean, answer matches choice)

\textbf{Decision:} Mark answer (D) $31$ confidently and move on
\end{verificationbox}

%\begin{roadmapbox}
\subsection*{A. Step-by-Step Solution Plan}

\textbf{Phase 1: Setup (30 seconds)}
\begin{enumerate}
    \item Read problem - logarithmic equation with variable bases
    \item Note $4 = 2^2$ and $8 = 2^3$
    \item Recognize need for change of base formula
    \item Prepare to use base 2
\end{enumerate}

\textbf{Phase 2: Change of Base Application (60 seconds)}
\begin{enumerate}
    \item Apply change of base to base 2:
    \[\frac{\log_2 4}{\log_2(3x)} = \frac{\log_2 8}{\log_2(2x)}\]
    \item Simplify numerators: $\frac{2}{\log_2(3x)} = \frac{3}{\log_2(2x)}$
    \item Cross-multiply: $3\log_2(3x) = 2\log_2(2x)$
\end{enumerate}

\textbf{Phase 3: Algebraic Solution (60 seconds)}
\begin{enumerate}
    \item Apply power rule: $\log_2[(3x)^3] = \log_2[(2x)^2]$
    \item Set arguments equal: $(3x)^3 = (2x)^2$
    \item Expand: $27x^3 = 4x^2$
    \item Divide by $x^2$ (valid since $x \neq 0$): $27x = 4$
    \item Solve: $x = \frac{4}{27}$
\end{enumerate}

\textbf{Phase 4: Final Answer (30 seconds)}
\begin{enumerate}
    \item Verify $\gcd(4,27) = 1$ (they share no common factors)
    \item Calculate $p + q = 4 + 27 = 31$
    \item Identify answer: (D) $31$
\end{enumerate}

\textbf{Total Time: 3 minutes}

\subsection*{B. Alternative Approaches}

\textbf{Approach 2: Power and Base Equalization}
\begin{enumerate}
    \item Rewrite: $\log_{(3x)^3} 4^3 = \log_{(2x)^2} 8^2$
    \item This gives: $\log_{(3x)^3} 64 = \log_{(2x)^2} 64$
    \item Therefore: $(3x)^3 = (2x)^2$
    \item Continue as before to get $x = \frac{4}{27}$
\end{enumerate}

\textbf{Approach 3: Exponential Form}
\begin{enumerate}
    \item Let $y = \log_{3x} 4 = \log_{2x} 8$
    \item Convert to exponential: $(3x)^y = 4$ and $(2x)^y = 8$
    \item Cube first equation: $(3x)^{3y} = 64$
    \item Square second equation: $(2x)^{2y} = 64$
    \item Therefore: $(3x)^{3y} = (2x)^{2y}$
    \item Taking $(1/y)$-th power: $(3x)^3 = (2x)^2$
    \item Solve as before
\end{enumerate}

\textbf{Approach 4: Natural Logarithm}
\begin{enumerate}
    \item Rewrite using $\ln$:
    \[2\log_{3x} 2 = 3\log_{2x} 2\]
    \item Convert to $\ln$: $\frac{2\ln 2}{\ln(3x)} = \frac{3\ln 2}{\ln(2x)}$
    \item Simplify: $\frac{2}{\ln(3x)} = \frac{3}{\ln(2x)}$
    \item Cross-multiply: $2\ln(2x) = 3\ln(3x)$
    \item Expand: $2\ln 2 + 2\ln x = 3\ln 3 + 3\ln x$
    \item Solve: $\ln x = 2\ln 2 - 3\ln 3 = \ln \frac{4}{27}$
    \item Therefore: $x = \frac{4}{27}$
\end{enumerate}

\subsection*{C. Checkpoints and Validation}

\begin{itemize}
    \item \textbf{Checkpoint 1}: After change of base, verify you have $\frac{2}{\log_2(3x)} = \frac{3}{\log_2(2x)}$
    \item \textbf{Checkpoint 2}: After cross-multiply, check $3\log_2(3x) = 2\log_2(2x)$
    \item \textbf{Checkpoint 3}: After power rule, verify $(3x)^3 = (2x)^2$
    \item \textbf{Checkpoint 4}: Final answer $\frac{4}{27}$ should satisfy domain restrictions
    \item \textbf{Checkpoint 5}: Answer $31$ should be among choices
\end{itemize}

\subsection*{D. Comprehensive Pitfall Guide}

\begin{highlightbox}{Common Errors to Avoid}
\textbf{Conceptual Errors:}
\begin{itemize}
    \item \textbf{Pitfall}: Forgetting to check domain restrictions
    \item \textbf{Prevention}: Always verify $3x > 0$, $3x \neq 1$, $2x > 0$, $2x \neq 1$
    
    \item \textbf{Pitfall}: Assuming $\log_{3x} 4 = \log_3 4 + \log_x 4$
    \item \textbf{Prevention}: Base is $(3x)$, not a product - use change of base correctly
    
    \item \textbf{Pitfall}: Canceling $x$ from both sides incorrectly
    \item \textbf{Prevention}: From $27x^3 = 4x^2$, divide by $x^2$, not $x^3$
    
    \item \textbf{Pitfall}: Not simplifying fraction to lowest terms
    \item \textbf{Prevention}: Check $\gcd(p,q) = 1$ before computing sum
\end{itemize}

\textbf{Calculation Errors:}
\begin{itemize}
    \item \textbf{Pitfall}: Computing $\log_2 4 = 4$ instead of $2$
    \item \textbf{Prevention}: $2^2 = 4$, so $\log_2 4 = 2$
    
    \item \textbf{Pitfall}: Computing $(3x)^3 = 9x^3$ instead of $27x^3$
    \item \textbf{Prevention}: $(3x)^3 = 3^3 \cdot x^3 = 27x^3$
    
    \item \textbf{Pitfall}: From $27x^3 = 4x^2$, getting $x = \frac{27}{4}$ instead of $\frac{4}{27}$
    \item \textbf{Prevention}: Divide both sides by $x^2$: $27x = 4$, so $x = \frac{4}{27}$
    
    \item \textbf{Pitfall}: Computing $4 + 27 = 30$ or $32$
    \item \textbf{Prevention}: Carefully: $4 + 27 = 31$
\end{itemize}

\textbf{Answer Choice Traps:}
\begin{itemize}
    \item \textbf{Pitfall}: Choosing (A) $5 = 1 + 4$ if missolving for $x = \frac{1}{4}$
    \item \textbf{Prevention}: Careful algebra in solving $27x^3 = 4x^2$
    
    \item \textbf{Pitfall}: Choosing (C) $17 = 8 + 9$ from incorrect simplification
    \item \textbf{Prevention}: Track exponents carefully through power rule
    
    \item \textbf{Pitfall}: Choosing (E) $35 = 8 + 27$ from $x = \frac{8}{27}$
    \item \textbf{Prevention}: Careful with $(2x)^2 = 4x^2$, not $2x^2$
\end{itemize}
\end{highlightbox}

%\end{roadmapbox}

\section*{Summary}

\textbf{Key Insight:} Logarithmic equations with variable bases require change of base formula. Recognizing $4 = 2^2$ and $8 = 2^3$ makes base 2 the natural choice.

\textbf{Answer: (D) $31$}

\textbf{Core Technique:} Change of base formula followed by power rule: 
\[\frac{\log_2 4}{\log_2(3x)} = \frac{\log_2 8}{\log_2(2x)} \Rightarrow (3x)^3 = (2x)^2 \Rightarrow x = \frac{4}{27}\]

\textbf{Time Required:} 3 minutes for experienced solver

\textbf{Difficulty Assessment:} Moderate (Problem 14/25) - requires solid logarithm knowledge but straightforward application

\begin{highlightbox}{Key Formulas Summary}
\textbf{Essential Logarithm Properties:}
\begin{itemize}
    \item \textbf{Change of Base}: $\log_b a = \frac{\log_c a}{\log_c b}$ for any valid base $c$
    \item \textbf{Power Rule}: $\log_b(a^n) = n\log_b a$
    \item \textbf{Product Rule}: $\log_b(xy) = \log_b x + \log_b y$
    \item \textbf{Quotient Rule}: $\log_b(x/y) = \log_b x - \log_b y$
    \item \textbf{Base Power Identity}: $\log_{b^n}(a^n) = \log_b a$
\end{itemize}

\textbf{For This Problem:}
\begin{align*}
\log_{3x} 4 &= \log_{2x} 8\\
\frac{2}{\log_2(3x)} &= \frac{3}{\log_2(2x)}\\
3\log_2(3x) &= 2\log_2(2x)\\
(3x)^3 &= (2x)^2\\
27x^3 &= 4x^2\\
x &= \frac{4}{27}\\
p + q &= 4 + 27 = 31
\end{align*}
\end{highlightbox}

\section*{Extended Analysis}

\subsection*{Why This Problem is Important}

\textbf{Mathematical Significance:}
\begin{itemize}
    \item \textbf{Change of Base Mastery}: Essential skill for all logarithm problems
    \item \textbf{Power Recognition}: Identifying $2^2$ and $2^3$ structure
    \item \textbf{Variable Base Handling}: More sophisticated than standard log equations
    \item \textbf{Domain Awareness}: Understanding when logarithms are defined
\end{itemize}

\textbf{Competition Value:}
\begin{itemize}
    \item \textbf{Standard Mid-Band Topic}: Tests core logarithm manipulation skills
    \item \textbf{Multiple Solution Paths}: Rewards both mechanical skill and insight
    \item \textbf{Clean Answer}: Rational solution with small numerator/denominator
    \item \textbf{Verification Possible}: Can check answer by substitution
\end{itemize}

\subsection*{Connection to Other Problems}

\textbf{Related AMC/AIME Topics:}
\begin{itemize}
    \item Logarithmic equations with numerical or algebraic bases
    \item Change of base formula applications
    \item Exponential equations and their logarithmic forms
    \item Systems involving logarithms and exponentials
    \item Domain restrictions in logarithmic functions
\end{itemize}

\textbf{Similar Problem Types:}
\begin{itemize}
    \item ``Solve $\log_a b = \log_b a$ for $a$ and $b$''
    \item ``Find all $x$ such that $\log_x 16 + \log_{2x} 4 = 3$''
    \item Logarithmic equations with reciprocal bases
    \item Problems involving $\log_b a \cdot \log_a b = 1$
\end{itemize}

\begin{competitionbox}
\subsection*{A. Time Management}

\textbf{Target Time: 3 minutes}

\textbf{Time Allocation:}
\begin{itemize}
    \item Reading and setup: 30 seconds
    \item Change of base application: 60 seconds
    \item Algebraic solution: 60 seconds
    \item Final answer and verification: 30 seconds
\end{itemize}

\textbf{Decision Points:}
\begin{itemize}
    \item If change of base seems complicated after 1 minute: try exponential form
    \item If stuck after 2 minutes: work backwards from answer choices
    \item If stuck after 4 minutes: make educated guess, flag for review
\end{itemize}

\subsection*{B. Problem Prioritization Strategy}

\textbf{Priority Level: High}

This is Problem 14, mid-band. It should be attempted after problems 1-13 are complete. The logarithm manipulation is standard for prepared students.

\textbf{Accessibility Factors:}
\begin{itemize}
    \item \textbf{Domain}: Algebra/Logarithms (if comfortable with logs, prioritize)
    \item \textbf{Structure}: Standard equation-solving format
    \item \textbf{Technique}: Change of base (core AMC technique)
    \item \textbf{Calculation}: Moderate algebra, clean answer
\end{itemize}

\subsection*{C. Risk Management}

\textbf{Confidence Threshold: 80\% to proceed}

\textbf{Risk Assessment:}
\begin{itemize}
    \item \textbf{High Risk}: Misapplying change of base formula
    \item \textbf{Medium Risk}: Algebraic errors in expansion
    \item \textbf{Low Risk}: Arithmetic in final sum (easily verified)
\end{itemize}

\textbf{Abandonment Criteria:}
\begin{itemize}
    \item If after 5 minutes you can't get $(3x)^3 = (2x)^2$, move on
    \item Educated guess: Try (D) $31$ (most natural-looking sum)
    \item Can return during review if time permits
\end{itemize}

\subsection*{D. Strategic Notes for Competition Day}

\begin{highlightbox}{Pro Tips}
\begin{itemize}
    \item \textbf{Pattern Recognition}: Variable bases in logarithms $\Rightarrow$ change of base formula
    \item \textbf{Power Recognition}: $4 = 2^2$ and $8 = 2^3$ suggests using base 2
    \item \textbf{Answer Form}: Answer must be sum of numerator and denominator in lowest terms
    \item \textbf{Domain Awareness}: Always check given restrictions make sense with your answer
    \item \textbf{Verification}: Can substitute back to check (time-consuming but reliable)
    \item \textbf{Answer Pattern}: (D) $31 = 4 + 27$ has natural factorization structure
\end{itemize}
\end{highlightbox}

\subsection*{Common Student Approaches}

\textbf{Approach Distribution (Estimated):}
\begin{itemize}
    \item 60\%: Change of base to base 2 (most common)
    \item 20\%: Change of base to natural log
    \item 10\%: Exponential form method
    \item 5\%: Power and base equalization
    \item 5\%: Other methods or trial-and-error
\end{itemize}

\textbf{Success Rates by Approach:}
\begin{itemize}
    \item Change of base (any base): 85\% success
    \item Exponential form: 75\% success
    \item Power equalization: 80\% success
    \item Trial from answer choices: 50\% success
\end{itemize}

\subsection*{Problem Design Analysis}

\textbf{Answer Choice Strategy:}
\begin{itemize}
    \item \textbf{(A) $5 = 1+4$}: Catches error if student gets $x = \frac{1}{4}$ or $x = 1$
    \item \textbf{(B) $13 = 4+9$}: Catches error if student gets $x = \frac{4}{9}$
    \item \textbf{(C) $17 = 8+9$}: Catches error if student gets $x = \frac{8}{9}$
    \item \textbf{(D) $31 = 4+27$}: Correct answer ($x = \frac{4}{27}$)
    \item \textbf{(E) $35 = 8+27$}: Catches error if student gets $x = \frac{8}{27}$
\end{itemize}

Each wrong answer corresponds to a specific algebraic mistake, making this a well-designed problem.

\subsection*{Teaching Points}

\textbf{For Students:}
\begin{enumerate}
    \item Always identify what base to use in change of base formula
    \item Look for power relationships ($4 = 2^2$, $8 = 2^3$)
    \item Check domain restrictions before and after solving
    \item Verify answer is in lowest terms before computing sum
    \item Multiple solution methods provide verification opportunities
\end{enumerate}

\textbf{For Instructors:}
\begin{enumerate}
    \item Emphasize change of base as fundamental tool
    \item Show multiple approaches to build flexibility
    \item Discuss why certain bases (like 2) are more natural
    \item Demonstrate verification by substitution
    \item Connect to exponential form for deeper understanding
\end{enumerate}

\subsection*{Extensions and Variations}

\textbf{Natural Extensions:}
\begin{enumerate}
    \item What if equation were $\log_{3x} 8 = \log_{2x} 4$?
    \item Solve $\log_{ax} b = \log_{cx} d$ for $x$ in terms of $a,b,c,d$
    \item Find all $x$ where $\log_{3x} 4 = \log_{2x} 8 = 2$
    \item What if bases were $3x$ and $3x^2$ instead of $3x$ and $2x$?
\end{enumerate}

\textbf{Harder Variations:}
\begin{enumerate}
    \item Solve $\log_{x+1} 9 = \log_{2x+1} 27$
    \item Find all $(x,y)$ with $\log_x y = \log_y x$ and $x + y = 10$
    \item Solve $\log_{3x} 4 + \log_{2x} 8 = 5$
    \item Find range of $x$ for which $\log_{3x} 4 < \log_{2x} 8$
\end{enumerate}

\section*{Logarithm Mastery Checklist}

\begin{highlightbox}{Essential Skills for This Problem}
$\square$ Know change of base formula cold\\
$\square$ Can recognize and use power rule instantly\\
$\square$ Understand domain restrictions for logarithms\\
$\square$ Can convert between logarithmic and exponential form\\
$\square$ Comfortable with algebraic manipulation\\
$\square$ Can verify GCD of two numbers\\
$\square$ Know when to use different logarithm bases
\end{highlightbox}

\begin{highlightbox}{Before Attempting This Problem}
$\square$ Have you reviewed change of base formula?\\
$\square$ Can you compute $\log_2 4$ and $\log_2 8$ mentally?\\
$\square$ Do you understand variable bases (like $3x$)?\\
$\square$ Can you expand $(3x)^3$ correctly?\\
$\square$ Have you allocated 3 minutes for this problem?\\
$\square$ Do you have scratch paper for algebra?
\end{highlightbox}

\begin{highlightbox}{During Problem Solving}
$\square$ Did you apply change of base correctly?\\
$\square$ Did you simplify $\log_2 4 = 2$ and $\log_2 8 = 3$?\\
$\square$ Did you cross-multiply properly?\\
$\square$ Did you apply power rule to get $(3x)^3 = (2x)^2$?\\
$\square$ Did you expand to $27x^3 = 4x^2$?\\
$\square$ Did you solve to get $x = \frac{4}{27}$?\\
$\square$ Did you verify $\gcd(4,27) = 1$?\\
$\square$ Did you compute $4 + 27 = 31$?
\end{highlightbox}

\begin{highlightbox}{After Solving}
$\square$ Confidence level $\geq 85\%$?\\
$\square$ Answer makes sense algebraically?\\
$\square$ Domain restrictions satisfied?\\
$\square$ Fraction in lowest terms?\\
$\square$ Answer is among choices?\\
$\square$ Ready to move to Problem 15?
\end{highlightbox}

\section*{Final Thoughts}

This problem exemplifies mid-band AMC 12 algebra: it requires knowing a key technique (change of base formula) and applying it systematically, but doesn't require deep insight or advanced methods. Success depends on:

\begin{itemize}
    \item \textbf{Technical proficiency} (knowing logarithm rules)
    \item \textbf{Pattern recognition} (seeing $4 = 2^2$, $8 = 2^3$)
    \item \textbf{Careful algebra} (expanding and solving correctly)
    \item \textbf{Attention to detail} (lowest terms, domain checks)
\end{itemize}

For students preparing for AMC 12, this problem serves as a benchmark: if you can solve it reliably in under 3 minutes, you have solid logarithm skills. If it takes longer or requires multiple attempts, review:
\begin{itemize}
    \item Change of base formula and its applications
    \item Power rule for logarithms
    \item Converting between logarithmic and exponential forms
    \item Domain restrictions for logarithmic functions
\end{itemize}

\subsection*{Key Takeaway}

The power of the change of base formula is that it reduces all logarithm comparisons to a common scale. In this problem:
\[\log_{3x} 4 = \log_{2x} 8 \quad\Rightarrow\quad \frac{\log_2 4}{\log_2(3x)} = \frac{\log_2 8}{\log_2(2x)}\]

This transforms an equation involving two different logarithms into a single algebraic equation in $x$. The ability to recognize when and how to apply this transformation is essential for AMC 12 success.

\vspace{1em}
\noindent\textbf{Final Answer: (D) 31}

\end{competitionbox}

\end{document}